% Distinguish original sampling from conditional empirical resampling.
\begin{tikzpicture}[
 x=1mm,y=1mm,
 every node/.style={font=\figurefont\fontsize{8}{10}\selectfont,text=FigureInk,inner sep=1pt},
 box/.style={draw=FigureBlue,line width=.7pt,rounded corners=1.8pt,minimum height=9mm,align=center,inner xsep=3pt},
 edge/.style={draw=FigureStroke,line width=.65pt,->},
 cond/.style={draw=FigureGreen,line width=.7pt,rounded corners=1.8pt,minimum height=9mm,align=center,inner xsep=3pt}
]
 \path[use as bounding box] (0,0) rectangle (145,57);
 \node[box,minimum width=16mm] (P) at (12,41) {总体 \(P\)};
 \node[box,minimum width=28mm] (Y) at (51,41) {原观测 \(\bm Y\)};
 \node[box,minimum width=24mm] (T) at (94,41) {估计 \(\widehat\tau\)};
 \draw[edge] (P) -- node[above=3pt] {独立抽样} (Y);
 \draw[edge] (Y) -- node[above=3pt] {同一统计量} (T);
 \node[anchor=west] at (113,41) {\(\widehat\tau-\tau(P)\)};
 \node[cond,minimum width=27mm] (E) at (51,14) {经验分布 \(\widehat P_n\)};
 \node[cond,minimum width=24mm] (YS) at (94,14) {重抽样 \(\bm Y^\ast\)};
 \node[cond,minimum width=18mm] (TS) at (132,14) {\(\widehat\tau^\ast\)};
 \draw[edge] (Y) -- node[right=3pt,align=left] {固定原观测\\记录等权 \(1/n\)} (E);
 \draw[edge] (E) -- node[above=3pt] {有放回抽样} (YS);
 \draw[edge] (YS) -- node[above=3pt] {重算} (TS);
 \node[anchor=west,text=FigureGreen] at (4,3) {给定 \(\bm Y\) 后，只平均重抽样随机性};
 \node[anchor=east] at (143,3) {\(\widehat\tau^\ast-\widehat\tau\)};
\end{tikzpicture}
